\section{CLS Templates ו-Reproducibility}

קובץ \texttt{.cls} (\eng{document class}) הוא לב העיצוב במסמכי \LaTeX{}.
המרצה הציג אותו כפרויקט עצמאי שמבטיח שכל מסמך שייצור הסוכן ייראה אחיד
--- עקרון הרפטביליות (\eng{Reproducibility}).

\subsection{מהו קובץ CLS?}

\eng{document class} מגדיר:
\begin{itemize}[rightmargin=2em]
  \item גודל עמוד, שוליים, פונטים (עברית ואנגלית).
  \item עיצוב כותרות, כותרת עליונה ותחתונה.
  \item פקודות מותאמות: \texttt{\textbackslash eng}, תיבות מידע, סגנון טבלאות.
  \item כללים קבועים: קופירייט, מספור, צבעים ארגוניים.
\end{itemize}

במקום שהסוכן יחזור על אותם הגדרות עיצוב בכל \texttt{main.tex}, הוא טוען:
\begin{lstlisting}
\documentclass{agentic-ai-template}
\end{lstlisting}

\subsection{CLS כ-Tool}

המרצה תיאר את ה-\eng{CLS} ככלי (\eng{tool}): ניתן להגדיר בו פונקציות
\LaTeX{} חדשות --- למשל נוסחאות במרכז שורה, בדיקת כיוון כתיבה, או
הוספת כותרת עליונה אוטומטית. כל תיקון עיצובי חוזר נכנס ל-\eng{CLS}
ולא נשאר ``תיקון חד-פעמי'' במסמך.

\subsection{לולאת תיקון רפטבילית}

כאשר מתגלה באג ב-\eng{PDF} (טבלה זזה ימינה, עברית הפוכה):
\begin{enumerate}[rightmargin=2em]
  \item מציגים לסוכן את הבעיה --- \textbf{בלי} לתקן מיד.
  \item מבקשים הסבר והצעת תיקון (בקרת איכות).
  \item אם הבעיה בשפה טבעית --- מעדכנים \eng{Skill}.
  \item אם הבעיה טכנית-עיצובית --- מעדכנים \eng{CLS}.
  \item מריצים קומפילציה מחדש; אם עבר --- הבעיה נפתרה לכל המסמכים העתידיים.
\end{enumerate}

\begin{warningbox}
תיקון נקודתי ב-\texttt{main.tex} במקום ב-\texttt{.cls} יגרום לבאג
לחזור במסמך הבא. זהו ניגוד ישיר לרפטביליות.
\end{warningbox}

\agenttable{
\caption{\eng{LaTeX}/סוכן מול \eng{Word} לרפטביליות}
\label{tab:reproducibility}
\begin{tabular}{@{}p{3.5cm}p{4.5cm}p{4.5cm}@{}}
\toprule
\textbf{מאפיין} & \textbf{LaTeX + CLS + Agent} & \textbf{Word + Agent} \\
\midrule
עקביות עיצוב & גבוהה (תבנית) & נמוכה \\
גרסאות & \eng{Git diff} ברור & קובץ בינארי \\
אוטומציה & קומפילציה מלאה & לעיתים ידני \\
עברית מורכבת & דורש \eng{LuaLaTeX} & בעיות נפוצות \\
\bottomrule
\end{tabular}
}

\subsection{יישום ארגוני}

בארגון, מומלץ:
\begin{itemize}[rightmargin=2em]
  \item תבנית \eng{CLS} אחת לכל סוג מסמך (דוח, הצעת מחיר, מאמר).
  \item \eng{PRD} קצר לכל פרויקט --- מגדיר מבנה וקהל יעד.
  \item סקריפט קומפילציה (\texttt{compile.ps1}) --- אותה פקודה לכל אחד.
\end{itemize}

\subsection{קשר למטלה 04}

מטלה זו \emph{מחייבת} קובץ \texttt{agentic-ai-template.cls} אמיתי.
ההערכה כוללת לא רק תוכן אלא גם את היכולת להפריד תוכן מעיצוב --- מיומנות
מרכזית בהגדרת הנחיות לסוכן.

\begin{insightbox}
``הסוכן שלכם יהיה רפטבילי --- הוא לא כל פעם ייצר משהו אחר'' (תמלול
ההרצאה). זהו עקרון ההנדסה שמאחורי מטלה 04.
\end{insightbox}
